% ============================================================
% Punkt 3: Implementacja i realizacja
% Projekt układu sumatora pełnego (Full Adder)
% z minimalizacją bramek logicznych
% ------------------------------------------------------------
% Plik przeznaczony do włączenia w dokumencie głównym Overleaf:
%   % ============================================================
% Punkt 3: Implementacja i realizacja
% Projekt układu sumatora pełnego (Full Adder)
% z minimalizacją bramek logicznych
% ------------------------------------------------------------
% Plik przeznaczony do włączenia w dokumencie głównym Overleaf:
%   % ============================================================
% Punkt 3: Implementacja i realizacja
% Projekt układu sumatora pełnego (Full Adder)
% z minimalizacją bramek logicznych
% ------------------------------------------------------------
% Plik przeznaczony do włączenia w dokumencie głównym Overleaf:
%   % ============================================================
% Punkt 3: Implementacja i realizacja
% Projekt układu sumatora pełnego (Full Adder)
% z minimalizacją bramek logicznych
% ------------------------------------------------------------
% Plik przeznaczony do włączenia w dokumencie głównym Overleaf:
%   \input{03_implementacja}
% Wymagane pakiety w dokumencie głównym:
%   \usepackage[utf8]{inputenc}  \usepackage[T1]{polski}
%   \usepackage{amsmath, amssymb}
% ============================================================

\section{Implementacja i realizacja}

\subsection{Wybór elementów logicznych}

Do realizacji projektu sumatora pełnego (Full Adder) wybrano podstawowe oraz złożone elementy cyfrowe z rodziny układów logicznych CMOS. Podstawowym kryterium wyboru była możliwość minimalizacji liczby struktur bramkowych oraz optymalizacja czasu propagacji sygnału. W projekcie wykorzystano:

\begin{itemize}
    \item Bramki XOR (alternatywy rozłącznej): stanowiące kluczowy element wyznaczający bit sumy. Bramka XOR generuje stan wysoki (1) na wyjściu tylko wtedy, gdy nieparzysta liczba sygnałów wejściowych jest w stanie wysokim.
    \item Bramki AND (koniunkcji): służące do utworzenia członów iloczynowych funkcji przeniesienia, tj. $a \cdot b$ oraz $c_1 \cdot c_{in}$.
    \item Bramkę OR (alternatywy): służącą do zsumowania logicznego członów iloczynowych, czego wynikiem jest sygnał przeniesienia $c_{out}$.
\end{itemize}

\subsection{Projekt układu sumatora}

Projekt obejmuje sumator pełny dodający trzy bity: $a$, $b$ oraz $c_{in}$. Układ posiada trzy wejścia binarne oraz dwa wyjścia cyfrowe: $s$, które przyjmuje wartość logiczną 1, gdy suma wejść jest nieparzysta, oraz $c_{out}$, które przyjmuje wartość logiczną 1, gdy $a + b + c_{in} \geq 2$.

Równania wyjściowe sumatora przed minimalizacją, zapisane w postaci sumy iloczynów (SOP) dla wszystkich stanów aktywnych wyjść, przybrałyby postać rozbudowanych funkcji logicznych składających się z czterech składników mintermów każda:
\begin{equation*}
s = \Sigma m(1, 2, 4, 7), \qquad c_{out} = \Sigma m(3, 5, 6, 7).
\end{equation*}

Zastosowanie bezpośredniego porównania par bitów za pomocą operacji XOR pozwala na zapisanie funkcji w postaci:
\begin{equation*}
s = a \oplus b \oplus c_{in}, \qquad c_{out} = a \cdot b \vee (a \oplus b) \cdot c_{in}.
\end{equation*}

Rozpisując bramkę XOR na podstawowe operacje logiczne (AND, OR, NOT) funkcja sumy przyjmuje postać:
\begin{equation*}
s = a \cdot b' \cdot c_{in}' \vee a \cdot b \cdot c_{in} \vee a' \cdot b \cdot c_{in}' \vee a' \cdot b' \cdot c_{in}.
\end{equation*}

Dzięki takiemu podejściu, zamiast budować rozległą siatkę bramek AND/OR dla ośmiu mintermów i implementować kilkanaście bramek, układ redukuje się do dwóch struktur sumujących oraz jednego elementu spinającego człony przeniesienia.

\subsection{Schemat logiczny układu}

Schemat logiczny zminimalizowanego sumatora opiera się na strukturze dwupoziomowej:

\begin{itemize}
    \item Poziom pierwszy (Porównanie bitów): bramka XOR pracująca nad parami bitów wejściowych. Bramka ta przyjmuje na wejścia bity $a$ i $b$, generując sygnał pośredni $c_1 = a \oplus b$.
    \item Poziom drugi (Wyznaczenie wyjść): sygnał $c_1$ podawany jest jednocześnie na drugą bramkę XOR (wraz z $c_{in}$), której wyjściem jest bit sumy $s = c_1 \oplus c_{in}$, oraz na bramkę AND (wraz z $c_{in}$), generującą człon $c_1 \cdot c_{in}$.
    \item Poziom trzeci (Agregacja wyniku): wyjście bramki AND oraz sygnał z drugiej bramki AND realizującej iloczyn $a \cdot b$ są doprowadzone do dwuwejściowej bramki OR. Sygnał wyjściowy jest bezpośrednim rezultatem operacji OR na sygnałach cząstkowych:
    \begin{equation*}
    c_{out} = a \cdot b \vee c_1 \cdot c_{in}.
    \end{equation*}
\end{itemize}

Przepływ sygnałów w układzie można przedstawić jako:
\begin{equation*}
a, b \xrightarrow{\text{XOR}} c_1 \xrightarrow{\text{XOR z } c_{in}} s, \qquad a \cdot b \xrightarrow{\text{OR}} c_{out} \xleftarrow{\text{OR}} c_1 \cdot c_{in}.
\end{equation*}

Kluczową zaletą tej struktury jest fakt, że sygnał pośredni $c_1$ wykorzystywany jest dwukrotnie -- zarówno w torze sumy, jak i w torze przeniesienia -- co eliminuje konieczność wykonywania tej samej operacji dwukrotnie.

Schemat logiczny zminimalizowanego układu przedstawia Rysunek~\ref{fig:schemat-fa}.

\begin{figure}[h]
    \centering
    \begin{circuitikz}[american]
        % --- bramka XOR 1 (a, b -> c1) ---
        \draw (0,4.2) node[xor port] (xor1) {};
        % --- bramka XOR 2 (c1, cin -> s) ---
        \draw (3.5,3.2) node[xor port] (xor2) {};
        % --- bramka AND 1 (a, b -> ab) ---
        \draw (0,1.6) node[and port] (and1) {};
        % --- bramka AND 2 (c1, cin -> c1cin) ---
        \draw (3.5,1.2) node[and port] (and2) {};
        % --- bramka OR (ab, c1cin -> cout) ---
        \draw (7,1.4) node[or port] (or1) {};
        % --- wejścia a, b ---
        \draw (xor1.in 1) -- ++(-1.2,0) node[left] {$a$} node[circ] {} coordinate (a);
        \draw (and1.in 1) -- (a) node[circ] {};
        \draw (xor1.in 2) -- ++(-1.2,0) node[left] {$b$} node[circ] {} coordinate (b);
        \draw (and1.in 2) -- (b) node[circ] {};
        % --- wejście cin (rozgałęzienie do xor2 i and2) ---
        \draw (xor2.in 2) -- ++(-1.2,0) coordinate (cin1) node[circ] {};
        \draw (and2.in 2) -- ++(-1.2,0) coordinate (cin2) node[circ] {};
        \draw (cin1) -- ++(-0.4,0) node[left] {$c_{in}$};
        \draw (cin1) -- (cin2);
        % --- sygnał c1: wyjście xor1 -> xor2.in1 i and2.in1 ---
        \draw (xor1.out) -- ++(0.35,0) node[circ] {} coordinate (c1);
        \draw (c1) |- (xor2.in 1);
        \draw (c1) |- (and2.in 1);
        % --- wyjście s ---
        \draw (xor2.out) -- ++(0.8,0) node[right] {$s$};
        % --- wyjście cout ---
        \draw (and1.out) -- ++(0.35,0) coordinate (ab) node[circ] {};
        \draw (ab) |- (or1.in 1);
        \draw (and2.out) -- ++(0.35,0) coordinate (cc) node[circ] {};
        \draw (cc) |- (or1.in 2);
        \draw (or1.out) -- ++(0.8,0) node[right] {$c_{out}$};
        % --- opisy sygnałów pośrednich ---
        \node[above] at (c1) {$c_1$};
        \node[above] at (1.9,2.5) {$a \cdot b$};
        \node[above] at (5.0,1.2) {$c_1 \cdot c_{in}$};
    \end{circuitikz}
    \caption{Schemat logiczny zminimalizowanego sumatora pełnego}
    \label{fig:schemat-fa}
\end{figure}

\subsection{Realizacja zminimalizowanego układu}

W praktycznej realizacji fizycznej lub strukturalnej (np. w strukturach reprogramowalnych FPGA lub na makiecie laboratoryjnej), zminimalizowany układ eliminuje potrzebę stosowania rozbudowanych siatek bramek AND/OR dla pojedynczych mintermów.

Przy użyciu standardowych układów scalonych TTL/CMOS, realizacja wymaga:

\begin{itemize}
    \item Jednego układu scalonego zawierającego bramki XOR (np. 74HC86 -- poczwórna bramka XOR 2-wejściowa).
    \item Jednego układu scalonego z bramkami AND (np. 74HC08 -- poczwórna bramka AND 2-wejściowa).
    \item Jednego układu scalonego z bramkami OR (np. 74HC32 -- poczwórna bramka OR 2-wejściowa).
\end{itemize}

Alternatywnie, jeżeli w dostępnym zestawie elementów nie występują bramki XOR, sumę można zrealizować wyłącznie na bramkach NAND, wykorzystując standardową tożsamość logiczną, co jednak zwiększa liczbę wykorzystanych struktur bramkowych.

Ograniczenie liczby połączeń do minimum znacząco redukuje ryzyko powstawania hazardu statycznego i dynamicznego w układzie, skracając jednocześnie łączny czas propagacji sygnału od wejścia do wyjścia do sumy czasów propagacji zaledwie trzech poziomów bramek.

% Wymagane pakiety w dokumencie głównym:
%   \usepackage[utf8]{inputenc}  \usepackage[T1]{polski}
%   \usepackage{amsmath, amssymb}
% ============================================================

\section{Implementacja i realizacja}

\subsection{Wybór elementów logicznych}

Do realizacji projektu sumatora pełnego (Full Adder) wybrano podstawowe oraz złożone elementy cyfrowe z rodziny układów logicznych CMOS. Podstawowym kryterium wyboru była możliwość minimalizacji liczby struktur bramkowych oraz optymalizacja czasu propagacji sygnału. W projekcie wykorzystano:

\begin{itemize}
    \item Bramki XOR (alternatywy rozłącznej): stanowiące kluczowy element wyznaczający bit sumy. Bramka XOR generuje stan wysoki (1) na wyjściu tylko wtedy, gdy nieparzysta liczba sygnałów wejściowych jest w stanie wysokim.
    \item Bramki AND (koniunkcji): służące do utworzenia członów iloczynowych funkcji przeniesienia, tj. $a \cdot b$ oraz $c_1 \cdot c_{in}$.
    \item Bramkę OR (alternatywy): służącą do zsumowania logicznego członów iloczynowych, czego wynikiem jest sygnał przeniesienia $c_{out}$.
\end{itemize}

\subsection{Projekt układu sumatora}

Projekt obejmuje sumator pełny dodający trzy bity: $a$, $b$ oraz $c_{in}$. Układ posiada trzy wejścia binarne oraz dwa wyjścia cyfrowe: $s$, które przyjmuje wartość logiczną 1, gdy suma wejść jest nieparzysta, oraz $c_{out}$, które przyjmuje wartość logiczną 1, gdy $a + b + c_{in} \geq 2$.

Równania wyjściowe sumatora przed minimalizacją, zapisane w postaci sumy iloczynów (SOP) dla wszystkich stanów aktywnych wyjść, przybrałyby postać rozbudowanych funkcji logicznych składających się z czterech składników mintermów każda:
\begin{equation*}
s = \Sigma m(1, 2, 4, 7), \qquad c_{out} = \Sigma m(3, 5, 6, 7).
\end{equation*}

Zastosowanie bezpośredniego porównania par bitów za pomocą operacji XOR pozwala na zapisanie funkcji w postaci:
\begin{equation*}
s = a \oplus b \oplus c_{in}, \qquad c_{out} = a \cdot b \vee (a \oplus b) \cdot c_{in}.
\end{equation*}

Rozpisując bramkę XOR na podstawowe operacje logiczne (AND, OR, NOT) funkcja sumy przyjmuje postać:
\begin{equation*}
s = a \cdot b' \cdot c_{in}' \vee a \cdot b \cdot c_{in} \vee a' \cdot b \cdot c_{in}' \vee a' \cdot b' \cdot c_{in}.
\end{equation*}

Dzięki takiemu podejściu, zamiast budować rozległą siatkę bramek AND/OR dla ośmiu mintermów i implementować kilkanaście bramek, układ redukuje się do dwóch struktur sumujących oraz jednego elementu spinającego człony przeniesienia.

\subsection{Schemat logiczny układu}

Schemat logiczny zminimalizowanego sumatora opiera się na strukturze dwupoziomowej:

\begin{itemize}
    \item Poziom pierwszy (Porównanie bitów): bramka XOR pracująca nad parami bitów wejściowych. Bramka ta przyjmuje na wejścia bity $a$ i $b$, generując sygnał pośredni $c_1 = a \oplus b$.
    \item Poziom drugi (Wyznaczenie wyjść): sygnał $c_1$ podawany jest jednocześnie na drugą bramkę XOR (wraz z $c_{in}$), której wyjściem jest bit sumy $s = c_1 \oplus c_{in}$, oraz na bramkę AND (wraz z $c_{in}$), generującą człon $c_1 \cdot c_{in}$.
    \item Poziom trzeci (Agregacja wyniku): wyjście bramki AND oraz sygnał z drugiej bramki AND realizującej iloczyn $a \cdot b$ są doprowadzone do dwuwejściowej bramki OR. Sygnał wyjściowy jest bezpośrednim rezultatem operacji OR na sygnałach cząstkowych:
    \begin{equation*}
    c_{out} = a \cdot b \vee c_1 \cdot c_{in}.
    \end{equation*}
\end{itemize}

Przepływ sygnałów w układzie można przedstawić jako:
\begin{equation*}
a, b \xrightarrow{\text{XOR}} c_1 \xrightarrow{\text{XOR z } c_{in}} s, \qquad a \cdot b \xrightarrow{\text{OR}} c_{out} \xleftarrow{\text{OR}} c_1 \cdot c_{in}.
\end{equation*}

Kluczową zaletą tej struktury jest fakt, że sygnał pośredni $c_1$ wykorzystywany jest dwukrotnie -- zarówno w torze sumy, jak i w torze przeniesienia -- co eliminuje konieczność wykonywania tej samej operacji dwukrotnie.

Schemat logiczny zminimalizowanego układu przedstawia Rysunek~\ref{fig:schemat-fa}.

\begin{figure}[h]
    \centering
    \begin{circuitikz}[american]
        % --- bramka XOR 1 (a, b -> c1) ---
        \draw (0,4.2) node[xor port] (xor1) {};
        % --- bramka XOR 2 (c1, cin -> s) ---
        \draw (3.5,3.2) node[xor port] (xor2) {};
        % --- bramka AND 1 (a, b -> ab) ---
        \draw (0,1.6) node[and port] (and1) {};
        % --- bramka AND 2 (c1, cin -> c1cin) ---
        \draw (3.5,1.2) node[and port] (and2) {};
        % --- bramka OR (ab, c1cin -> cout) ---
        \draw (7,1.4) node[or port] (or1) {};
        % --- wejścia a, b ---
        \draw (xor1.in 1) -- ++(-1.2,0) node[left] {$a$} node[circ] {} coordinate (a);
        \draw (and1.in 1) -- (a) node[circ] {};
        \draw (xor1.in 2) -- ++(-1.2,0) node[left] {$b$} node[circ] {} coordinate (b);
        \draw (and1.in 2) -- (b) node[circ] {};
        % --- wejście cin (rozgałęzienie do xor2 i and2) ---
        \draw (xor2.in 2) -- ++(-1.2,0) coordinate (cin1) node[circ] {};
        \draw (and2.in 2) -- ++(-1.2,0) coordinate (cin2) node[circ] {};
        \draw (cin1) -- ++(-0.4,0) node[left] {$c_{in}$};
        \draw (cin1) -- (cin2);
        % --- sygnał c1: wyjście xor1 -> xor2.in1 i and2.in1 ---
        \draw (xor1.out) -- ++(0.35,0) node[circ] {} coordinate (c1);
        \draw (c1) |- (xor2.in 1);
        \draw (c1) |- (and2.in 1);
        % --- wyjście s ---
        \draw (xor2.out) -- ++(0.8,0) node[right] {$s$};
        % --- wyjście cout ---
        \draw (and1.out) -- ++(0.35,0) coordinate (ab) node[circ] {};
        \draw (ab) |- (or1.in 1);
        \draw (and2.out) -- ++(0.35,0) coordinate (cc) node[circ] {};
        \draw (cc) |- (or1.in 2);
        \draw (or1.out) -- ++(0.8,0) node[right] {$c_{out}$};
        % --- opisy sygnałów pośrednich ---
        \node[above] at (c1) {$c_1$};
        \node[above] at (1.9,2.5) {$a \cdot b$};
        \node[above] at (5.0,1.2) {$c_1 \cdot c_{in}$};
    \end{circuitikz}
    \caption{Schemat logiczny zminimalizowanego sumatora pełnego}
    \label{fig:schemat-fa}
\end{figure}

\subsection{Realizacja zminimalizowanego układu}

W praktycznej realizacji fizycznej lub strukturalnej (np. w strukturach reprogramowalnych FPGA lub na makiecie laboratoryjnej), zminimalizowany układ eliminuje potrzebę stosowania rozbudowanych siatek bramek AND/OR dla pojedynczych mintermów.

Przy użyciu standardowych układów scalonych TTL/CMOS, realizacja wymaga:

\begin{itemize}
    \item Jednego układu scalonego zawierającego bramki XOR (np. 74HC86 -- poczwórna bramka XOR 2-wejściowa).
    \item Jednego układu scalonego z bramkami AND (np. 74HC08 -- poczwórna bramka AND 2-wejściowa).
    \item Jednego układu scalonego z bramkami OR (np. 74HC32 -- poczwórna bramka OR 2-wejściowa).
\end{itemize}

Alternatywnie, jeżeli w dostępnym zestawie elementów nie występują bramki XOR, sumę można zrealizować wyłącznie na bramkach NAND, wykorzystując standardową tożsamość logiczną, co jednak zwiększa liczbę wykorzystanych struktur bramkowych.

Ograniczenie liczby połączeń do minimum znacząco redukuje ryzyko powstawania hazardu statycznego i dynamicznego w układzie, skracając jednocześnie łączny czas propagacji sygnału od wejścia do wyjścia do sumy czasów propagacji zaledwie trzech poziomów bramek.

% Wymagane pakiety w dokumencie głównym:
%   \usepackage[utf8]{inputenc}  \usepackage[T1]{polski}
%   \usepackage{amsmath, amssymb}
% ============================================================

\section{Implementacja i realizacja}

\subsection{Wybór elementów logicznych}

Do realizacji projektu sumatora pełnego (Full Adder) wybrano podstawowe oraz złożone elementy cyfrowe z rodziny układów logicznych CMOS. Podstawowym kryterium wyboru była możliwość minimalizacji liczby struktur bramkowych oraz optymalizacja czasu propagacji sygnału. W projekcie wykorzystano:

\begin{itemize}
    \item Bramki XOR (alternatywy rozłącznej): stanowiące kluczowy element wyznaczający bit sumy. Bramka XOR generuje stan wysoki (1) na wyjściu tylko wtedy, gdy nieparzysta liczba sygnałów wejściowych jest w stanie wysokim.
    \item Bramki AND (koniunkcji): służące do utworzenia członów iloczynowych funkcji przeniesienia, tj. $a \cdot b$ oraz $c_1 \cdot c_{in}$.
    \item Bramkę OR (alternatywy): służącą do zsumowania logicznego członów iloczynowych, czego wynikiem jest sygnał przeniesienia $c_{out}$.
\end{itemize}

\subsection{Projekt układu sumatora}

Projekt obejmuje sumator pełny dodający trzy bity: $a$, $b$ oraz $c_{in}$. Układ posiada trzy wejścia binarne oraz dwa wyjścia cyfrowe: $s$, które przyjmuje wartość logiczną 1, gdy suma wejść jest nieparzysta, oraz $c_{out}$, które przyjmuje wartość logiczną 1, gdy $a + b + c_{in} \geq 2$.

Równania wyjściowe sumatora przed minimalizacją, zapisane w postaci sumy iloczynów (SOP) dla wszystkich stanów aktywnych wyjść, przybrałyby postać rozbudowanych funkcji logicznych składających się z czterech składników mintermów każda:
\begin{equation*}
s = \Sigma m(1, 2, 4, 7), \qquad c_{out} = \Sigma m(3, 5, 6, 7).
\end{equation*}

Zastosowanie bezpośredniego porównania par bitów za pomocą operacji XOR pozwala na zapisanie funkcji w postaci:
\begin{equation*}
s = a \oplus b \oplus c_{in}, \qquad c_{out} = a \cdot b \vee (a \oplus b) \cdot c_{in}.
\end{equation*}

Rozpisując bramkę XOR na podstawowe operacje logiczne (AND, OR, NOT) funkcja sumy przyjmuje postać:
\begin{equation*}
s = a \cdot b' \cdot c_{in}' \vee a \cdot b \cdot c_{in} \vee a' \cdot b \cdot c_{in}' \vee a' \cdot b' \cdot c_{in}.
\end{equation*}

Dzięki takiemu podejściu, zamiast budować rozległą siatkę bramek AND/OR dla ośmiu mintermów i implementować kilkanaście bramek, układ redukuje się do dwóch struktur sumujących oraz jednego elementu spinającego człony przeniesienia.

\subsection{Schemat logiczny układu}

Schemat logiczny zminimalizowanego sumatora opiera się na strukturze dwupoziomowej:

\begin{itemize}
    \item Poziom pierwszy (Porównanie bitów): bramka XOR pracująca nad parami bitów wejściowych. Bramka ta przyjmuje na wejścia bity $a$ i $b$, generując sygnał pośredni $c_1 = a \oplus b$.
    \item Poziom drugi (Wyznaczenie wyjść): sygnał $c_1$ podawany jest jednocześnie na drugą bramkę XOR (wraz z $c_{in}$), której wyjściem jest bit sumy $s = c_1 \oplus c_{in}$, oraz na bramkę AND (wraz z $c_{in}$), generującą człon $c_1 \cdot c_{in}$.
    \item Poziom trzeci (Agregacja wyniku): wyjście bramki AND oraz sygnał z drugiej bramki AND realizującej iloczyn $a \cdot b$ są doprowadzone do dwuwejściowej bramki OR. Sygnał wyjściowy jest bezpośrednim rezultatem operacji OR na sygnałach cząstkowych:
    \begin{equation*}
    c_{out} = a \cdot b \vee c_1 \cdot c_{in}.
    \end{equation*}
\end{itemize}

Przepływ sygnałów w układzie można przedstawić jako:
\begin{equation*}
a, b \xrightarrow{\text{XOR}} c_1 \xrightarrow{\text{XOR z } c_{in}} s, \qquad a \cdot b \xrightarrow{\text{OR}} c_{out} \xleftarrow{\text{OR}} c_1 \cdot c_{in}.
\end{equation*}

Kluczową zaletą tej struktury jest fakt, że sygnał pośredni $c_1$ wykorzystywany jest dwukrotnie -- zarówno w torze sumy, jak i w torze przeniesienia -- co eliminuje konieczność wykonywania tej samej operacji dwukrotnie.

Schemat logiczny zminimalizowanego układu przedstawia Rysunek~\ref{fig:schemat-fa}.

\begin{figure}[h]
    \centering
    \begin{circuitikz}[american]
        % --- bramka XOR 1 (a, b -> c1) ---
        \draw (0,4.2) node[xor port] (xor1) {};
        % --- bramka XOR 2 (c1, cin -> s) ---
        \draw (3.5,3.2) node[xor port] (xor2) {};
        % --- bramka AND 1 (a, b -> ab) ---
        \draw (0,1.6) node[and port] (and1) {};
        % --- bramka AND 2 (c1, cin -> c1cin) ---
        \draw (3.5,1.2) node[and port] (and2) {};
        % --- bramka OR (ab, c1cin -> cout) ---
        \draw (7,1.4) node[or port] (or1) {};
        % --- wejścia a, b ---
        \draw (xor1.in 1) -- ++(-1.2,0) node[left] {$a$} node[circ] {} coordinate (a);
        \draw (and1.in 1) -- (a) node[circ] {};
        \draw (xor1.in 2) -- ++(-1.2,0) node[left] {$b$} node[circ] {} coordinate (b);
        \draw (and1.in 2) -- (b) node[circ] {};
        % --- wejście cin (rozgałęzienie do xor2 i and2) ---
        \draw (xor2.in 2) -- ++(-1.2,0) coordinate (cin1) node[circ] {};
        \draw (and2.in 2) -- ++(-1.2,0) coordinate (cin2) node[circ] {};
        \draw (cin1) -- ++(-0.4,0) node[left] {$c_{in}$};
        \draw (cin1) -- (cin2);
        % --- sygnał c1: wyjście xor1 -> xor2.in1 i and2.in1 ---
        \draw (xor1.out) -- ++(0.35,0) node[circ] {} coordinate (c1);
        \draw (c1) |- (xor2.in 1);
        \draw (c1) |- (and2.in 1);
        % --- wyjście s ---
        \draw (xor2.out) -- ++(0.8,0) node[right] {$s$};
        % --- wyjście cout ---
        \draw (and1.out) -- ++(0.35,0) coordinate (ab) node[circ] {};
        \draw (ab) |- (or1.in 1);
        \draw (and2.out) -- ++(0.35,0) coordinate (cc) node[circ] {};
        \draw (cc) |- (or1.in 2);
        \draw (or1.out) -- ++(0.8,0) node[right] {$c_{out}$};
        % --- opisy sygnałów pośrednich ---
        \node[above] at (c1) {$c_1$};
        \node[above] at (1.9,2.5) {$a \cdot b$};
        \node[above] at (5.0,1.2) {$c_1 \cdot c_{in}$};
    \end{circuitikz}
    \caption{Schemat logiczny zminimalizowanego sumatora pełnego}
    \label{fig:schemat-fa}
\end{figure}

\subsection{Realizacja zminimalizowanego układu}

W praktycznej realizacji fizycznej lub strukturalnej (np. w strukturach reprogramowalnych FPGA lub na makiecie laboratoryjnej), zminimalizowany układ eliminuje potrzebę stosowania rozbudowanych siatek bramek AND/OR dla pojedynczych mintermów.

Przy użyciu standardowych układów scalonych TTL/CMOS, realizacja wymaga:

\begin{itemize}
    \item Jednego układu scalonego zawierającego bramki XOR (np. 74HC86 -- poczwórna bramka XOR 2-wejściowa).
    \item Jednego układu scalonego z bramkami AND (np. 74HC08 -- poczwórna bramka AND 2-wejściowa).
    \item Jednego układu scalonego z bramkami OR (np. 74HC32 -- poczwórna bramka OR 2-wejściowa).
\end{itemize}

Alternatywnie, jeżeli w dostępnym zestawie elementów nie występują bramki XOR, sumę można zrealizować wyłącznie na bramkach NAND, wykorzystując standardową tożsamość logiczną, co jednak zwiększa liczbę wykorzystanych struktur bramkowych.

Ograniczenie liczby połączeń do minimum znacząco redukuje ryzyko powstawania hazardu statycznego i dynamicznego w układzie, skracając jednocześnie łączny czas propagacji sygnału od wejścia do wyjścia do sumy czasów propagacji zaledwie trzech poziomów bramek.

% Wymagane pakiety w dokumencie głównym:
%   \usepackage[utf8]{inputenc}  \usepackage[T1]{polski}
%   \usepackage{amsmath, amssymb}
% ============================================================

\section{Implementacja i realizacja}

\subsection{Wybór elementów logicznych}

Do realizacji projektu sumatora pełnego (Full Adder) wybrano podstawowe oraz złożone elementy cyfrowe z rodziny układów logicznych CMOS. Podstawowym kryterium wyboru była możliwość minimalizacji liczby struktur bramkowych oraz optymalizacja czasu propagacji sygnału. W projekcie wykorzystano:

\begin{itemize}
    \item Bramki XOR (alternatywy rozłącznej): stanowiące kluczowy element wyznaczający bit sumy. Bramka XOR generuje stan wysoki (1) na wyjściu tylko wtedy, gdy nieparzysta liczba sygnałów wejściowych jest w stanie wysokim.
    \item Bramki AND (koniunkcji): służące do utworzenia członów iloczynowych funkcji przeniesienia, tj. $a \cdot b$ oraz $c_1 \cdot c_{in}$.
    \item Bramkę OR (alternatywy): służącą do zsumowania logicznego członów iloczynowych, czego wynikiem jest sygnał przeniesienia $c_{out}$.
\end{itemize}

\subsection{Projekt układu sumatora}

Projekt obejmuje sumator pełny dodający trzy bity: $a$, $b$ oraz $c_{in}$. Układ posiada trzy wejścia binarne oraz dwa wyjścia cyfrowe: $s$, które przyjmuje wartość logiczną 1, gdy suma wejść jest nieparzysta, oraz $c_{out}$, które przyjmuje wartość logiczną 1, gdy $a + b + c_{in} \geq 2$.

Równania wyjściowe sumatora przed minimalizacją, zapisane w postaci sumy iloczynów (SOP) dla wszystkich stanów aktywnych wyjść, przybrałyby postać rozbudowanych funkcji logicznych składających się z czterech składników mintermów każda:
\begin{equation*}
s = \Sigma m(1, 2, 4, 7), \qquad c_{out} = \Sigma m(3, 5, 6, 7).
\end{equation*}

Zastosowanie bezpośredniego porównania par bitów za pomocą operacji XOR pozwala na zapisanie funkcji w postaci:
\begin{equation*}
s = a \oplus b \oplus c_{in}, \qquad c_{out} = a \cdot b \vee (a \oplus b) \cdot c_{in}.
\end{equation*}

Rozpisując bramkę XOR na podstawowe operacje logiczne (AND, OR, NOT) funkcja sumy przyjmuje postać:
\begin{equation*}
s = a \cdot b' \cdot c_{in}' \vee a \cdot b \cdot c_{in} \vee a' \cdot b \cdot c_{in}' \vee a' \cdot b' \cdot c_{in}.
\end{equation*}

Dzięki takiemu podejściu, zamiast budować rozległą siatkę bramek AND/OR dla ośmiu mintermów i implementować kilkanaście bramek, układ redukuje się do dwóch struktur sumujących oraz jednego elementu spinającego człony przeniesienia.

\subsection{Schemat logiczny układu}

Schemat logiczny zminimalizowanego sumatora opiera się na strukturze dwupoziomowej:

\begin{itemize}
    \item Poziom pierwszy (Porównanie bitów): bramka XOR pracująca nad parami bitów wejściowych. Bramka ta przyjmuje na wejścia bity $a$ i $b$, generując sygnał pośredni $c_1 = a \oplus b$.
    \item Poziom drugi (Wyznaczenie wyjść): sygnał $c_1$ podawany jest jednocześnie na drugą bramkę XOR (wraz z $c_{in}$), której wyjściem jest bit sumy $s = c_1 \oplus c_{in}$, oraz na bramkę AND (wraz z $c_{in}$), generującą człon $c_1 \cdot c_{in}$.
    \item Poziom trzeci (Agregacja wyniku): wyjście bramki AND oraz sygnał z drugiej bramki AND realizującej iloczyn $a \cdot b$ są doprowadzone do dwuwejściowej bramki OR. Sygnał wyjściowy jest bezpośrednim rezultatem operacji OR na sygnałach cząstkowych:
    \begin{equation*}
    c_{out} = a \cdot b \vee c_1 \cdot c_{in}.
    \end{equation*}
\end{itemize}

Przepływ sygnałów w układzie można przedstawić jako:
\begin{equation*}
a, b \xrightarrow{\text{XOR}} c_1 \xrightarrow{\text{XOR z } c_{in}} s, \qquad a \cdot b \xrightarrow{\text{OR}} c_{out} \xleftarrow{\text{OR}} c_1 \cdot c_{in}.
\end{equation*}

Kluczową zaletą tej struktury jest fakt, że sygnał pośredni $c_1$ wykorzystywany jest dwukrotnie -- zarówno w torze sumy, jak i w torze przeniesienia -- co eliminuje konieczność wykonywania tej samej operacji dwukrotnie.

Schemat logiczny zminimalizowanego układu przedstawia Rysunek~\ref{fig:schemat-fa}.

\begin{figure}[h]
    \centering
    \begin{circuitikz}[american]
        % --- bramka XOR 1 (a, b -> c1) ---
        \draw (0,4.2) node[xor port] (xor1) {};
        % --- bramka XOR 2 (c1, cin -> s) ---
        \draw (3.5,3.2) node[xor port] (xor2) {};
        % --- bramka AND 1 (a, b -> ab) ---
        \draw (0,1.6) node[and port] (and1) {};
        % --- bramka AND 2 (c1, cin -> c1cin) ---
        \draw (3.5,1.2) node[and port] (and2) {};
        % --- bramka OR (ab, c1cin -> cout) ---
        \draw (7,1.4) node[or port] (or1) {};
        % --- wejścia a, b ---
        \draw (xor1.in 1) -- ++(-1.2,0) node[left] {$a$} node[circ] {} coordinate (a);
        \draw (and1.in 1) -- (a) node[circ] {};
        \draw (xor1.in 2) -- ++(-1.2,0) node[left] {$b$} node[circ] {} coordinate (b);
        \draw (and1.in 2) -- (b) node[circ] {};
        % --- wejście cin (rozgałęzienie do xor2 i and2) ---
        \draw (xor2.in 2) -- ++(-1.2,0) coordinate (cin1) node[circ] {};
        \draw (and2.in 2) -- ++(-1.2,0) coordinate (cin2) node[circ] {};
        \draw (cin1) -- ++(-0.4,0) node[left] {$c_{in}$};
        \draw (cin1) -- (cin2);
        % --- sygnał c1: wyjście xor1 -> xor2.in1 i and2.in1 ---
        \draw (xor1.out) -- ++(0.35,0) node[circ] {} coordinate (c1);
        \draw (c1) |- (xor2.in 1);
        \draw (c1) |- (and2.in 1);
        % --- wyjście s ---
        \draw (xor2.out) -- ++(0.8,0) node[right] {$s$};
        % --- wyjście cout ---
        \draw (and1.out) -- ++(0.35,0) coordinate (ab) node[circ] {};
        \draw (ab) |- (or1.in 1);
        \draw (and2.out) -- ++(0.35,0) coordinate (cc) node[circ] {};
        \draw (cc) |- (or1.in 2);
        \draw (or1.out) -- ++(0.8,0) node[right] {$c_{out}$};
        % --- opisy sygnałów pośrednich ---
        \node[above] at (c1) {$c_1$};
        \node[above] at (1.9,2.5) {$a \cdot b$};
        \node[above] at (5.0,1.2) {$c_1 \cdot c_{in}$};
    \end{circuitikz}
    \caption{Schemat logiczny zminimalizowanego sumatora pełnego}
    \label{fig:schemat-fa}
\end{figure}

\subsection{Realizacja zminimalizowanego układu}

W praktycznej realizacji fizycznej lub strukturalnej (np. w strukturach reprogramowalnych FPGA lub na makiecie laboratoryjnej), zminimalizowany układ eliminuje potrzebę stosowania rozbudowanych siatek bramek AND/OR dla pojedynczych mintermów.

Przy użyciu standardowych układów scalonych TTL/CMOS, realizacja wymaga:

\begin{itemize}
    \item Jednego układu scalonego zawierającego bramki XOR (np. 74HC86 -- poczwórna bramka XOR 2-wejściowa).
    \item Jednego układu scalonego z bramkami AND (np. 74HC08 -- poczwórna bramka AND 2-wejściowa).
    \item Jednego układu scalonego z bramkami OR (np. 74HC32 -- poczwórna bramka OR 2-wejściowa).
\end{itemize}

Alternatywnie, jeżeli w dostępnym zestawie elementów nie występują bramki XOR, sumę można zrealizować wyłącznie na bramkach NAND, wykorzystując standardową tożsamość logiczną, co jednak zwiększa liczbę wykorzystanych struktur bramkowych.

Ograniczenie liczby połączeń do minimum znacząco redukuje ryzyko powstawania hazardu statycznego i dynamicznego w układzie, skracając jednocześnie łączny czas propagacji sygnału od wejścia do wyjścia do sumy czasów propagacji zaledwie trzech poziomów bramek.
